\section{Composite-Score Analysis}

\subsection{Component and head ablations}

Table~\ref{tab:scoreablation} does not support a claim that the composite score is universally superior. Coordinator uncertainty achieves a slightly larger mean AUROC, while the composite score provides a 0.0196 mean recall gain. The paired differences are not significant after Holm correction. The defensible conclusion is that the composite construction provides a recall-oriented balanced operating score whose individual signals have interpretable roles. Anomaly-only and residual-only scores are insufficient.

\begin{table}[H]
\caption{Open-set score ablation with the same 24-byte message, five seeds, and eight holdouts.}
\label{tab:scoreablation}
\centering
\begin{tabular}{lrrrr}
\toprule
Open-set score & Mean AUROC & Min AUROC & Mean recall & Min recall \\
\midrule
Composite $S_{\rm COS}$ & 0.9022 & 0.4401 & 0.7238 & 0.0000 \\
Coordinator uncertainty & 0.9050 & 0.3838 & 0.7042 & 0.0000 \\
Uncertainty+anomaly & 0.9014 & 0.4563 & 0.7056 & 0.0000 \\
Anomaly only & 0.8296 & 0.5524 & 0.0270 & 0.0000 \\
Message residual only & 0.7711 & 0.3099 & 0.0000 & 0.0000 \\
\bottomrule
\end{tabular}
\end{table}

\subsection{Hyperparameter sensitivity}

The reference weights were fixed after the exploratory development stage and before the five-seed confirmatory rerun. The full $5^3$ grid uses $\{0,0.25,0.5,0.75,1\}$ for each parameter. The selected point obtains mean AUROC 0.9022 and mean recall 0.7238. The largest mean AUROC is 0.9045 at $\lambda=0.75,\gamma=0$ (independent of $\beta$), with recall 0.7148. The largest mean recall is 0.7281 at $(\beta,\lambda,\gamma)=(0.25,0.50,0.75)$, with AUROC 0.8960. Thus, the selected score lies within 0.0023 AUROC of the best ranking point and within 0.0043 recall of the best recall point. No grid point solves UDPFlood: the minimum family recall remains zero.

\begin{table}[H]
\caption{Representative composite-score settings over 40 matched trials.}
\label{tab:sensitivity}
\centering
\begin{tabular}{lrrrrrr}
\toprule
Setting & $\beta$ & $\lambda$ & $\gamma$ & Mean AUROC & Mean recall & Max FPR \\
\midrule
Reference compromise & 0.25 & 0.50 & 0.25 & 0.9022 & 0.7238 & 0.0518 \\
Highest mean AUROC & 0.00 & 0.75 & 0.00 & 0.9045 & 0.7148 & 0.0514 \\
Highest mean recall & 0.25 & 0.50 & 0.75 & 0.8960 & 0.7281 & 0.0512 \\
Anomaly branch only & 0.00 & 0.00 & 0.00 & 0.8296 & 0.0667 & 0.0513 \\
\bottomrule
\end{tabular}
\end{table}

\begin{figure}[H]
\centering
\includegraphics[width=0.82\textwidth]{figures/sensitivity_auroc.pdf}
\caption{Sensitivity of mean AUROC to the composite-score weights. The selected setting belongs to a broad high-performing region rather than a sharp isolated optimum.}
\label{fig:sensitivity}
\end{figure}

\subsection{Paired statistical comparisons}

Table~\ref{tab:statistics} reports selected paired comparisons for the 24-byte float32 reference; Section~6.3 establishes its detection equivalence to 16Q. Relative to the all-source logit ensemble, X-MAG-COS improves AUROC and recall significantly after Holm correction. Relative to the centralized Random Forest, the open-set differences are not significant. Relative to \dhmethod{}, neither the small AUROC decrease nor the small recall increase is significant. These results limit the supported claim to a communication-efficient operating point rather than universal statistical superiority.

\begin{table}[H]
\caption{Selected paired tests over 40 seed--holdout pairs. Confidence intervals are bootstrap 95\% intervals for X-MAG-COS-24B minus the comparator; $p_{\rm W,H}$ is the Holm-adjusted Wilcoxon p-value.}
\label{tab:statistics}
\centering
\scriptsize
\begin{tabular}{L{0.25\textwidth}lrrr}
\toprule
Comparator & Metric & Mean difference [95\% CI] & $p_{\rm W,H}$ \\
\midrule
\dhmethod{} & AUROC & $-0.0028$ [$-0.0082$, $0.0034$] & 0.162 \\
\dhmethod{} & Recall & $0.0196$ [$-0.0026$, $0.0409$] & 0.224 \\
All-source logit average & AUROC & $0.0517$ [$0.0305$, $0.0736$] & $<0.001$ \\
All-source logit average & Recall & $0.2134$ [$0.0834$, $0.3421$] & 0.015 \\
Central Random Forest & AUROC & $0.0149$ [$-0.0413$, $0.0757$] & 1.000 \\
Central Random Forest & Recall & $-0.0748$ [$-0.1499$, $-0.0058$] & 1.000 \\
FedAvg linear & Known F1 & $0.0974$ [$0.0801$, $0.1142$] & $<0.001$ \\
Central Extra Trees & Known F1 & $0.0046$ [$0.0025$, $0.0071$] & $<0.001$ \\
\bottomrule
\end{tabular}
\end{table}

\section{Attribution-Proxy Fidelity and Cost}

The TreeSHAP comparison is deliberately limited to the two difficult held-out families and seed 42. Table~\ref{tab:shap} shows high absolute-rank Spearman correlation but only moderate exact top-1 agreement. The proxy therefore preserves broad feature ordering better than exact per-flow feature identity. This evidence supports the term ``attribution-aware'' but not an equivalence claim with SHAP or LIME.

\begin{table}[H]
\caption{Attribution-proxy agreement with TreeSHAP. Top-3 overlap is the mean Jaccard similarity.}
\label{tab:shap}
\centering
\begin{tabular}{lrrrrr}
\toprule
Holdout & Samples & Top-1 agreement & Top-3 overlap & Rank $\rho$ & Speedup \\
\midrule
UDPFlood & 1,896 & 0.271 & 0.650 & 0.937 & $156.7\times$ \\
SlowrateDoS & 1,851 & 0.409 & 0.416 & 0.908 & $73.1\times$ \\
\bottomrule
\end{tabular}
\end{table}

The proxy requires approximately 0.028 s for each sampled hard-holdout set, whereas TreeSHAP requires 2.03--4.33 s. This cost difference motivates the proxy for per-flow evidence generation. The modest top-1 agreement remains a limitation: the transmitted field is a ranking surrogate, not a complete local explanation.

\section{Channel Artifacts and Targeted Evidence Manipulation}

\subsection{Exploratory peer-context simulation}

To move beyond independent deterministic partitions, we evaluated an optional coordinator input that concatenates each owner message with the pooled latest delivered messages from the other source monitors. Packet loss and bounded event-order jitter emulate channel artifacts. Table~\ref{tab:network} shows that naive peer pooling is not robust: performance changes non-monotonically and often collapses under small loss/jitter. This negative result prevents us from presenting peer-context fusion as a validated contribution.

\begin{table}[H]
\caption{Exploratory pooled-peer-context simulation at seed 42. Values use the conformal 5\% rule. The non-monotonic results indicate unstable score calibration under naive asynchronous pooling.}
\label{tab:network}
\centering
\small
\begin{tabular}{llrrrr}
\toprule
Holdout & Channel & Delivery & Known F1 & AUROC & Unknown recall \\
\midrule
\multirow{4}{*}{UDPFlood}
& Ideal & 1.000 & 0.9986 & 0.6763 & 0.0000 \\
& Edge (1\%, 3 steps) & 0.990 & 0.9838 & 0.3005 & 0.0102 \\
& Congested (5\%, 20 steps) & 0.950 & 0.9436 & 0.1346 & 0.0500 \\
& Severe (10\%, 50 steps) & 0.900 & 0.9516 & 0.8041 & 0.0000 \\
\midrule
\multirow{4}{*}{SlowrateDoS}
& Ideal & 1.000 & 0.9990 & 0.9873 & 1.0000 \\
& Edge (1\%, 3 steps) & 0.990 & 0.9270 & 0.2198 & 0.0136 \\
& Congested (5\%, 20 steps) & 0.951 & 0.9312 & 0.3165 & 0.0503 \\
& Severe (10\%, 50 steps) & 0.901 & 0.9020 & 0.4980 & 0.0000 \\
\bottomrule
\end{tabular}
\end{table}

\subsection{Active source compromise}

The targeted experiment compromises one active source monitor in each hard-holdout task. Because the ownership metadata occupies fewer active slots than the nominal eight, requested fractions of 10--30\% all selected the same one source; the study therefore supports a single-source-compromise conclusion, not a dose-response claim.

Class-evidence suppression, benign-message replay, and combined manipulation collapse known macro-F1 to 0.0966 for UDPFlood and 0.0737 for SlowrateDoS. Attribution replacement preserves known F1 but substantially changes rejection. Anomaly suppression reduces unknown recall to zero for UDPFlood and 0.0154 for SlowrateDoS. The method is consequently not robust to an authenticated source being malicious. Deployment requires message authentication, source attestation, reliability weighting, or redundant observation.

\begin{table}[H]
\caption{Targeted single-source compromise at seed 42. FPR is reported because some attacks obtain apparent recall by exceeding the 5\% known-traffic budget.}
\label{tab:active}
\centering
\scriptsize
\begin{tabular}{llrrrr}
\toprule
Holdout & Manipulation & Known F1 & AUROC & Recall & FPR \\
\midrule
\multirow{5}{*}{UDPFlood}
& Class suppression & 0.0966 & 0.6064 & 0.0000 & 0.0000 \\
& Attribution replacement & 0.9986 & 0.4822 & 0.0000 & 0.0462 \\
& Anomaly suppression & 0.9987 & 0.4664 & 0.0000 & 0.1490 \\
& Benign-message replay & 0.0966 & 0.6258 & 0.0000 & 0.0000 \\
& Combined & 0.0966 & 0.6296 & 0.0000 & 0.0000 \\
\midrule
\multirow{5}{*}{SlowrateDoS}
& Class suppression & 0.0737 & 0.9142 & 0.0033 & 0.0226 \\
& Attribution replacement & 0.9990 & 0.9238 & 0.5031 & 0.0820 \\
& Anomaly suppression & 0.9990 & 0.8855 & 0.0154 & 0.0208 \\
& Benign-message replay & 0.0737 & 0.9725 & 0.4819 & 0.0109 \\
& Combined & 0.0737 & 0.9160 & 0.0000 & 0.0075 \\
\bottomrule
\end{tabular}
\end{table}
